\chapter{सांख्यिकी और प्रायिकता}\label{ch:g9:statproba}

आँकड़े और संयोग हर जगह हैं: खेलों के नतीजे, मौसम की भविष्यवाणी, पासे
और ताश।
यह अध्याय सिखाता है कि संख्याओं की किसी श्रेणी को उसके \omterm{def:g9:statproba:indicators}{माध्य},
\omterm{def:g9:statproba:indicators}{माध्यिका} और \omterm{def:g9:statproba:indicators}{परास} में कैसे समेटा जाए, और सादे यादृच्छिक प्रयोगों की
\omterm{def:g9:statproba:probability}{प्रायिकता} कैसे निकाली जाए --- दो चरणों वाले प्रयोग भी, जिन्हें वृक्ष
आरेख सँभाल लेते हैं। दोनों विषय इस श्रृंखला के हाई-स्कूल खंड में
कहीं आगे तक ले जाए गए हैं।

\section{आँकड़ों को समेटना}

\begin{definition}[माध्य, माध्यिका, परास]\label{def:g9:statproba:indicators}
$N$ संख्याओं की किसी श्रेणी के लिए:
\begin{itemize}
  \item \emph{माध्य}\index{माध्य} सारे मानों का \omterm{def:g1:addition:def}{योगफल} है, $N$ से भाग
        दिया हुआ;
  \item \emph{माध्यिका}\index{माध्यिका} वह मान है जो क्रम से लगी
        श्रेणी को बराबर आकार के दो हिस्सों में बाँट देता है: $N$
        विषम हो तो बीच वाला मान, और $N$ सम हो तो बीच वाले दो मानों
        का मध्यबिंदु;
  \item \emph{परास}\index{परास} सबसे बड़े मान में से सबसे छोटा मान
        घटाने पर मिलता है।
\end{itemize}
\end{definition}

\begin{example}\label{ex:g9:statproba:indicators}
किसी विद्यार्थी के अंक: $8,\ 12,\ 9,\ 15,\ 11$।
\begin{enumerate}
  \item \omterm{def:g9:statproba:indicators}{माध्य}: $\dfrac{8 + 12 + 9 + 15 + 11}{5} = \dfrac{55}{5} = 11$।
  \item \omterm{def:g9:statproba:indicators}{माध्यिका}: श्रेणी को क्रम से लगाओ: $8, 9, 11, 12, 15$; बीच
        वाला (तीसरा) मान $11$ है।
  \item \omterm{def:g9:statproba:indicators}{परास}: $15 - 8 = 7$।
\end{enumerate}
$17$ का छठा अंक जुड़ जाए: क्रम से $8, 9, 11, 12, 15, 17$, तो
\omterm{def:g9:statproba:indicators}{माध्यिका} $11$ और $12$ का मध्यबिंदु यानी $11.5$ हो जाती है, और \omterm{def:g9:statproba:indicators}{माध्य}
$\frac{72}{6} = 12$।
\end{example}

\begin{example}[भारित माध्य]\label{ex:g9:statproba:weighted}
एक दिन में बिके जूतों के नाप:
\begin{center}
\renewcommand{\arraystretch}{1.2}
\begin{tabular}{c|ccccc}
नाप & $38$ & $39$ & $40$ & $41$ & $42$ \\
\hline
\omterm{def:g7:stats:frequency}{गिनती} & $3$ & $7$ & $6$ & $3$ & $1$ \\
\end{tabular}
\end{center}
नाप का \omterm{def:g9:statproba:indicators}{माध्य} है
\[
\frac{3 \times 38 + 7 \times 39 + 6 \times 40 + 3 \times 41 + 1 \times 42}
{3 + 7 + 6 + 3 + 1}
= \frac{114 + 273 + 240 + 123 + 42}{20}
= \frac{792}{20} = 39.6 .
\]
\end{example}

\begin{remark}
\omterm{def:g9:statproba:indicators}{माध्य} और \omterm{def:g9:statproba:indicators}{माध्यिका} में बहुत फ़र्क़ पड़ सकता है। श्रेणी
$10, 10, 10, 10, 60$ में \omterm{def:g9:statproba:indicators}{माध्य} $20$ है (मान $60$ ने उसे ऊपर खींच
लिया) जबकि \omterm{def:g9:statproba:indicators}{माध्यिका} $10$ है। हमेशा पूछो कि कौन-सा सूचक हालत को बेहतर
बताता है।
\end{remark}

\section{प्रायिकता}

\begin{definition}[प्रायिकता]\label{def:g9:statproba:probability}
किसी यादृच्छिक प्रयोग के कई संभव \emph{परिणाम} होते हैं; और
\emph{घटना} परिणामों का एक समुच्चय है। हर परिणाम को एक
\emph{प्रायिकता}\index{प्रायिकता} मिलती है: $0$ और $1$ के बीच की एक
संख्या, और सारे परिणामों की प्रायिकताओं का योग $1$ होता है; किसी
घटना $A$ की प्रायिकता $\P(A)$ उसके परिणामों की प्रायिकताओं का \omterm{def:g1:addition:def}{योगफल}
है। जब सभी $n$ परिणाम समान रूप से संभावित हों, तब
\[
\P(A) = \frac{\text{अनुकूल परिणामों की गिनती}}{n}.
\]
\end{definition}

\begin{example}\label{ex:g9:statproba:spinner}
कोई चक्र $8$ बराबर खंडों में बँटा है: $3$ लाल, $3$ नीले, $2$ हरे। हर
खंड की \omterm{def:g9:statproba:probability}{प्रायिकता} $\frac18$ है, इसलिए
\[
\P(\text{लाल}) = \frac38, \qquad
\P(\text{हरा}) = \frac28 = \frac14, \qquad
\P(\text{हरा नहीं}) = 1 - \frac14 = \frac34 .
\]
\end{example}

\begin{proposition}[पूरक]\label{prop:g9:statproba:complement}
हर घटना $A$ के लिए उसकी \emph{पूरक घटना} $\bar A$ (``$A$ घटित नहीं
होती'') यह शर्त मानती है:
\[
\P(\bar A) = 1 - \P(A).
\]
\end{proposition}

\begin{proof}
हर परिणाम $A$ और $\bar A$ में से ठीक एक में पड़ता है, इसलिए
$\P(A) + \P(\bar A)$ सारे परिणामों की प्रायिकताओं को जोड़ देता है,
जो $1$ है।
\end{proof}

\section{दो चरणों वाले प्रयोग}

\begin{method}[वृक्ष आरेख]\label{met:g9:statproba:tree}
दो चरणों वाले किसी प्रयोग के लिए:
\begin{enumerate}
  \item पहले चरण के हर परिणाम के लिए एक शाखा खींचो, फिर हर शाखा को
        दूसरे चरण के परिणामों से आगे बढ़ाओ, और हर शाखा पर उसकी
        \omterm{def:g9:statproba:probability}{प्रायिकता} लिख दो;
  \item किसी रास्ते की \omterm{def:g9:statproba:probability}{प्रायिकता} पाने के लिए उस रास्ते भर की
        प्रायिकताओं को गुणा करो;
  \item किसी घटना को बनाने वाले सारे रास्तों की प्रायिकताएँ जोड़ दो।
\end{enumerate}
\end{method}

\begin{example}\label{ex:g9:statproba:tree}
किसी थैली में $2$ लाल और $3$ काली गोटियाँ हैं। एक गोटी निकालो, उसका
रंग नोट करो, \emph{उसे वापस डाल दो}, और फिर से निकालो। हर बार लाल
$\frac25$ \omterm{def:g9:statproba:probability}{प्रायिकता} से आती है और काली $\frac35$ \omterm{def:g9:statproba:probability}{प्रायिकता} से।

\begin{omfigure}
\begin{tikzpicture}[scale=0.95]
  \coordinate (root) at (0,0);
  \coordinate (R) at (2.2,1.0);
  \coordinate (B) at (2.2,-1.0);
  \coordinate (RR) at (4.7,1.55);
  \coordinate (RB) at (4.7,0.45);
  \coordinate (BR) at (4.7,-0.45);
  \coordinate (BB) at (4.7,-1.55);
  \draw (root) -- (R) node[midway, above, font=\small, sloped] {$\frac25$};
  \draw (root) -- (B) node[midway, below, font=\small, sloped] {$\frac35$};
  \draw (R) -- (RR) node[midway, above, font=\small, sloped] {$\frac25$};
  \draw (R) -- (RB) node[midway, below, font=\small, sloped] {$\frac35$};
  \draw (B) -- (BR) node[midway, above, font=\small, sloped] {$\frac25$};
  \draw (B) -- (BB) node[midway, below, font=\small, sloped] {$\frac35$};
  \node[omThm, font=\small, above left=-2pt] at (R) {ल};
  \node[font=\small, below left=-2pt] at (B) {क};
  \node[right, font=\small] at (RR)
    {लल:\ $\frac25 \times \frac25 = \frac{4}{25}$};
  \node[right, font=\small] at (RB)
    {लक:\ $\frac25 \times \frac35 = \frac{6}{25}$};
  \node[right, font=\small] at (BR)
    {कल:\ $\frac35 \times \frac25 = \frac{6}{25}$};
  \node[right, font=\small] at (BB)
    {कक:\ $\frac35 \times \frac35 = \frac{9}{25}$};
\end{tikzpicture}

{\small वापस डालकर किए गए दोनों निकालों का वृक्ष: शाखाओं के साथ-साथ
गुणा करो, और जाँच लो कि चारों रास्तों का योग $1$ है।}
\end{omfigure}

एक ही रंग की दो गोटियाँ आने की \omterm{def:g9:statproba:probability}{प्रायिकता}:
$\frac{4}{25} + \frac{9}{25} = \frac{13}{25}$। कम से कम एक लाल आने
की \omterm{def:g9:statproba:probability}{प्रायिकता}: $1 - \P(\text{कक}) = 1 - \frac{9}{25} = \frac{16}{25}$।
\end{example}

\begin{remark}[आवृत्तियाँ प्रायिकताओं के पास पहुँचती हैं]
किसी न्यायसंगत पासे को $6000$ बार फेंकने पर \emph{क़रीब} $1000$ छक्के
आते हैं --- ठीक-ठीक नहीं। दोहराव जितने बढ़ते जाते हैं, देखी गई
आवृत्तियाँ प्रायिकताओं के उतने ही पास आती जाती हैं: इसी से \omterm{def:g9:statproba:probability}{प्रायिकता}
बार-बार दोहराए जाने वाले प्रयोगों का सही औज़ार बन जाती है (यह कहानी
इस श्रृंखला के हाई-स्कूल खंड में आगे बढ़ती है)।
\end{remark}

\section{अभ्यास}

\begin{exercise}[$\star$]\label{exo:g9:statproba:1}
श्रेणी $14,\ 9,\ 17,\ 9,\ 12,\ 11,\ 12$ का \omterm{def:g9:statproba:indicators}{माध्य}, \omterm{def:g9:statproba:indicators}{माध्यिका} और \omterm{def:g9:statproba:indicators}{परास}
निकालो।
\end{exercise}

\begin{exercise}[$\star$]\label{exo:g9:statproba:2}
हफ़्ते भर दोपहर के तापमान $19,\ 21,\ 24,\ 18,\ 22,\ 25,\ 19$ डिग्री
रहे। \omterm{def:g9:statproba:indicators}{माध्य} (दसवें भाग तक) और \omterm{def:g9:statproba:indicators}{माध्यिका} निकालो।
\end{exercise}

\begin{exercise}[$\star$]\label{exo:g9:statproba:3}
किसी टीम ने $10$ मुक़ाबलों में जितने गोल किए:
\begin{center}
\renewcommand{\arraystretch}{1.2}
\begin{tabular}{c|cccc}
गोल & $0$ & $1$ & $2$ & $3$ \\
\hline
मुक़ाबले & $2$ & $4$ & $3$ & $1$ \\
\end{tabular}
\end{center}
हर मुक़ाबले में गोलों का \omterm{def:g9:statproba:indicators}{माध्य} निकालो, और \omterm{def:g9:statproba:indicators}{माध्यिका} भी।
\end{exercise}

\begin{exercise}[$\star$]\label{exo:g9:statproba:4}
कोई न्यायसंगत पासा एक बार फेंका जाता है। इनकी \omterm{def:g9:statproba:probability}{प्रायिकता} निकालो:
``$3$ आना''; ``\omterm{def:g3:numbers:evenodd}{सम संख्या} आना''; ``कम से कम $3$ आना''; ``$6$ न
आना''।
\end{exercise}

\begin{exercise}[$\star$]\label{exo:g9:statproba:5}
किसी डिब्बे में $12$ गेंदें हैं: $5$ सफ़ेद, $4$ लाल और $3$ हरी; एक
गेंद यादृच्छिक ढंग से निकाली जाती है। $\P(\text{सफ़ेद})$,
$\P(\text{लाल या हरी})$ और $\P(\text{सफ़ेद नहीं})$ निकालो। आख़िरी दो
के बारे में क्या दिखाई देता है?
\end{exercise}

\begin{exercise}[$\star\star$]\label{exo:g9:statproba:6}
एक सिक्का उछाला जाता है, फिर एक न्यायसंगत पासा फेंका जाता है। वृक्ष
खींचो (या परिणाम गिनो) और चित \emph{और} $6$ आने की \omterm{def:g9:statproba:probability}{प्रायिकता} निकालो;
फिर चित \emph{या} $6$ (या दोनों) आने की।
\end{exercise}

\begin{exercise}[$\star\star$]\label{exo:g9:statproba:7}
\cref{ex:g9:statproba:tree} वाली थैली में ($2$ लाल, $3$ काली) अब
दोनों गोटियाँ \emph{बिना} वापस डाले निकाली जाती हैं। नया वृक्ष खींचो
(ध्यान रहे: दूसरे स्तर की प्रायिकताएँ बदल जाती हैं) और दो काली
गोटियाँ आने की \omterm{def:g9:statproba:probability}{प्रायिकता} निकालो, फिर अलग-अलग रंग की दो गोटियाँ आने
की।
\end{exercise}

\begin{exercise}[$\star\star$]\label{exo:g9:statproba:8}
$9$ मुक़ाबलों के बाद किसी बास्केटबॉल खिलाड़ी का \omterm{def:g9:statproba:indicators}{माध्य} $14$ अंक प्रति
मुक़ाबला है।
\begin{enumerate}
  \item उसने कुल कितने अंक बनाए हैं?
  \item \omterm{def:g9:statproba:indicators}{माध्य} $15$ करने के लिए उसे दसवें मुक़ाबले में कितने अंक
        बनाने होंगे?
\end{enumerate}
\end{exercise}

\begin{exercise}[$\star\star\star$]\label{exo:g9:statproba:9}
एक खेल: दो न्यायसंगत पासे फेंको; दोनों के नतीजे बराबर आ जाएँ
(``जोड़ा''), तो तुम जीते।
\begin{enumerate}
  \item जीतने की \omterm{def:g9:statproba:probability}{प्रायिकता} निकालो।
  \item तुम दो बार खेलते हो (दोनों खेल एक-दूसरे से स्वतंत्र)। वृक्ष
        की मदद से कम से कम एक बार जीतने की \omterm{def:g9:statproba:probability}{प्रायिकता} निकालो।
\end{enumerate}
\end{exercise}

\section{समस्या: शेवालिए की डुबो देने वाली बाज़ियाँ}

\begin{problem}[{सप्ताहांत समस्या --- $1654$ का वह जुआ-झगड़ा जिसने
प्रायिकता का सिद्धांत खड़ा कर दिया, और वह जन्मदिन वाला संयोग जो हर
कक्षा को चकमा दे जाता है}]
\label{pb:g9:statproba:1}
$1654$ में जुए के शौक़ीन एक फ़्रांसीसी रईस, शेवालिए दे मेरे, पासे की
एक बाज़ी से मालामाल हो गए और एक दूसरी बाज़ी में हारने लगे, जिसे वे
पहली के बराबर मानते थे। उलझन में उन्होंने गणितज्ञ ब्लेज़ पास्कल से
पूछा; पास्कल ने पिएर दे फ़र्मा को लिखा; और उन दोनों की चिट्ठियों ने
\omterm{def:g9:statproba:probability}{प्रायिकता} के सिद्धांत की नींव रख दी। यह समस्या पूरा क़िस्सा इसी
अध्याय के औज़ारों से दोहराती है --- समान रूप से संभावित परिणाम, वृक्ष
और पूरक का नियम (\cref{prop:g9:statproba:complement}) --- और सबसे
अनहोनी लगने वाले संयोग पर जाकर ख़त्म होती है।

\medskip
\noindent\textbf{भाग I --- पासे, ईमानदारी से गिने हुए।}
\begin{enumerate}
  \item कोई न्यायसंगत पासा दो बार फेंको। वृक्ष (या $6 \times 6$ की
        सारणी) की मदद से दोनों बार छक्का \emph{न} आने की \omterm{def:g9:statproba:probability}{प्रायिकता}
        निकालो, और उससे कम से कम एक छक्का आने की \omterm{def:g9:statproba:probability}{प्रायिकता}।
  \item कोई सहपाठी कहता है: ``हर बार छह में से एक मौक़ा, और दो बार,
        यानी $\frac16 + \frac16 = \frac13$।'' सवाल 1 से तुलना करो
        और ठीक-ठीक बताओ कि उसका जोड़ किन परिणामों को दो बार गिन
        जाता है।
  \item दो पासे फेंककर जोड़ना: \omterm{def:g1:addition:def}{योगफल} $7$ आने की और \omterm{def:g1:addition:def}{योगफल} $2$ आने की
        \omterm{def:g9:statproba:probability}{प्रायिकता} निकालो। $7$ जुआरियों का मनपसंद क्यों है?
  \item पुराना झगड़ा: दो पासों से \omterm{def:g1:addition:def}{योगफल} $9$ ज़्यादा संभव है या
        \omterm{def:g1:addition:def}{योगफल} $10$? परिणाम गिनकर फ़ैसला करो।
  \item सवाल 1 को आम बनाओ: समझाओ कि $n$ बार फेंकने पर कम से कम एक
        छक्का आने की \omterm{def:g9:statproba:probability}{प्रायिकता}
        \[
        1 - \left(\frac56\right)^{\!n} ,
        \]
        क्यों होती है --- वृक्ष की ``छक्का नहीं'' वाली शाखाएँ एक
        स्तर से अगले स्तर तक गुणा होती चली जाती हैं।
\end{enumerate}

\medskip
\noindent\textbf{भाग II --- शेवालिए की दोनों बाज़ियाँ।}
\begin{enumerate}[resume]
  \item पहली बाज़ी, बराबर दाँव पर: \emph{पासे की चार फेंकों में कम
        से कम एक छक्का}। $\left(\frac56\right)^4$ को \omterm{def:g9:fractions:fraction}{भिन्न} के रूप
        में और दशमलव में निकालो, और शेवालिए के जीतने की \omterm{def:g9:statproba:probability}{प्रायिकता}
        भी। क्या बाज़ी अच्छी थी?
  \item शेवालिए का अपना तर्क यह था: ``चार फेंकें, हर एक
        $\frac16$ की: यानी $\frac46$ का मौक़ा।'' यहाँ उससे जवाब
        लगभग सही निकल आया --- पर उसी तर्क को $7$ फेंकों तक बढ़ाओ:
        कौन-सी बेतुकी बात निकलती है, और सवाल 2 वाली कौन-सी ग़लती
        वह दोहरा रहा है?
  \item दूसरी बाज़ी, जिसे वे \omterm{def:g7:prop:table}{समानुपात} के भरोसे पहली के बराबर मानते
        थे (``$24$ फेंकें, हर एक $\frac1{36}$ की: फिर वही
        $\frac{24}{36} = \frac46$''): \emph{दो पासों की $24$
        फेंकों में कम से कम एक बार दोहरा छक्का}। जीतने की असली
        \omterm{def:g9:statproba:probability}{प्रायिकता} $1 - \left(\frac{35}{36}\right)^{24}$ निकालो
        (कैलकुलेटर से)। शेवालिए धीरे-धीरे क्यों डूब गए?
  \item दो पासों की कितनी फेंकें उनका पलड़ा फिर से भारी कर देतीं?
        कैलकुलेटर से $n = 25$ आज़माओ और नतीजा निकालो।
  \item दो वाक्यों में सीख लिखो: किसी \omterm{def:g9:statproba:probability}{प्रायिकता} को कोशिशों की \omterm{def:g7:stats:frequency}{गिनती}
        से गुणा कर देने में आम तौर पर ख़राबी क्या है --- और उस
        लालच की जगह कौन-सा सही औज़ार आता है?
\end{enumerate}

\medskip
\noindent\textbf{भाग III --- \omterm{def:g9:statproba:indicators}{माध्य}, माध्यिकाएँ, जन्मदिन।}
\begin{enumerate}[resume]
  \item कोई छोटी कंपनी अपने नौ कर्मचारियों को महीने के $2\,000$
        यूरो देती है और निदेशक को $20\,000$। वेतन का \omterm{def:g9:statproba:indicators}{माध्य} और
        \omterm{def:g9:statproba:indicators}{माध्यिका} निकालो (\cref{def:g9:statproba:indicators})।
        नौकरी के विज्ञापन में ईमानदारी से कौन-सी संख्या छापी जानी
        चाहिए?
  \item पाँच अंकों की ऐसी सूची बनाओ जिसका \omterm{def:g9:statproba:indicators}{माध्य} $12$ हो और
        \omterm{def:g9:statproba:indicators}{माध्यिका} $15$। यह जोड़ी (\omterm{def:g9:statproba:indicators}{माध्य}, \omterm{def:g9:statproba:indicators}{माध्यिका} से नीचे) अंकों की
        शक्ल के बारे में क्या कहती है?
  \item जन्मदिन वाली समस्या: $23$ विद्यार्थियों की किसी कक्षा में
        कम से कम दो का जन्मदिन एक ही दिन पड़ने की \omterm{def:g9:statproba:probability}{प्रायिकता} क्या
        है? पूरक घटना खड़ी करो (तेईसों जन्मदिन अलग-अलग):
        \[
        \frac{365}{365} \times \frac{364}{365} \times
        \frac{363}{365} \times \dots \times \frac{343}{365} .
        \]
        पहले तीन गुणकों का \omterm{def:g2:mult:def}{गुणनफल} निकालो, बताओ कि गुणक आगे कैसे
        बदलते जाते हैं, और --- यह मानकर कि पूरा \omterm{def:g2:mult:def}{गुणनफल} क़रीब
        $0.493$ है --- सवाल का जवाब दो। ज़्यादातर लोग ``बहुत
        मुश्किल'' का अंदाज़ा लगाते हैं: गणित क्या कहता है?
  \item अटकल की मरम्मत: $23$ विद्यार्थियों की कक्षा में कितने
        \emph{जोड़े} एक ही जन्मदिन बाँट सकते हैं
        (\cref{pb:g6:wholes:1} वाली हाथ मिलाने की \omterm{def:g7:stats:frequency}{गिनती})? एक वाक्य
        में समझाओ कि हैरानी की जड़ यही संख्या क्यों है, $23$ ख़ुद
        क्यों नहीं।
  \item समापन: दोनों में से किसी एक नतीजे को जाँचने का कोई प्रयोग
        बताओ --- असली पासे से दे मेरे की बाज़ी, या कई कक्षाओं में
        जन्मदिन वाली समस्या --- और ध्यान से लिखो कि दोहराव बढ़ने
        के साथ देखी गई आवृत्तियों से तुम्हें क्या उम्मीद है। (उस
        उम्मीद का एक नाम है, बृहत् संख्याओं का नियम; उसका प्रमाण
        हाई-स्कूल खंड में इंतज़ार कर रहा है।)

\end{enumerate}
\end{problem}
